\section{Broadcasts and Shoutcasting}
\label{sec:literature-review:broadcast-and-shoutcasting}

\subsection{Broadcasting Medium Tools}
\label{subsec:broadcast-medium-tools}

In the previous chapter, Section \ref{sec:introduction:esports_and_shoutcasting}, the increasing broadcasting and audience engagement in esports events were introduced. As a result of this growth, several studies have been conducted in order to improve the understanding of this domain, and explore possible applications of Artificial Intelligence to enhance the viewer experience. 

\medskip

A key advantage for research in esports broadcasts is the availability of large-scale data. The broadcasting of esports events is often done through online streaming platforms, such as Twitch and YouTube, which provide a large amount of data in the form of video recordings (VODs) and live chat interactions. Multiple studies include the use of the broadcasting sites as a source of data for analyzing audience engagement. For instance, Pedrassoli et al. \cite{pedrassoli2024passive} provided an interactive tool for the platform Twitch, which showed live game statistics, while also obtaining telemetry data from the audience through the entirety of a Dota 2 tournament.

\medskip

The telemetry data obtained from the audience provided insights to understand the moments in which users tend to be more engaged with the broadcast. It also quantified user engagement, by defining an \textit{active user time} metric of around 5 seconds, when the game is in a critical moment.

\medskip

The importance of the active user time metric comes from the insight that it provides about the short-term user engagement during critical game events. These findings can aid in the design of temporal segments correlated to discourse functions while modeling shoutcasting structure. However, a limitation from this is that it heavily relies on the nature of the game. While the present study also focuses on a game from the MOBA genre, the same values of active user time may not be applicable without refinement according to the specific pacing of \textit{League of Legends}.

\medskip

Similar studies have also been conducted in order to find more subtle insights, through less structured data, such as the content of chat interactions. The research conducted by Jubaer et al. \cite{jubaer2024analyzing} analyzed the live data of the chat interactions of the official YouTube channels of PUBG Mobile in Bangladesh, India, and Pakistan, through the use of pre-trained language models for sentiment analysis and topic discovery. 

\medskip

Apart from discovering the positive sentiment of the audience, which aligns with the increasing popularity of esports, the study also found a correlation between game events / outcomes, and the sentiments of the audience, while topics were often a variety of:

\begin{itemize}
  \item Strategic Analysis
  \item Player Performance
  \item Community Interactions
\end{itemize}

\medskip

Two main insights were obtained from the study. The first one is that there is a notable overlap between the topics of discussion in the live chat and the contents of the shoutcasting. This means that some of these topics could serve as a guide for the design of discourse function labels. The second one is that the models used show resilience when working with live chat data, which is often noisy and unstructured. This means that the use of pre-trained language models could be a good option for the design of an AI system to analyze the discourse of shoutcasters, as it would be able to handle the complexity of the language used in esports broadcasts.

\medskip

A limitation of this study is that it focuses only on the official YouTube channels of PUBG Mobile in a specific regional context. In contrast, the broadcasts of League of Legends are available in both their official esports channel, and licensed third-party channels. This results in a more varied dataset, as it will allow to observe different styles of shoutcasting and production, which may be more representative of the overall shoutcasting discourse, but also more difficult to model. 

\medskip

Overall, these studies show the potential of obtaining data directly from broadcasting platforms, as they provide a large amount of unstructured data. However, these are mostly focused on the audience engagement metrics, and not on the analysis of the shoutcasting discourse which faciliates the engagement. This is a gap that the present study aims to fill, by directly modeling the structure and function of shoutcasting discourse, through some of the same tools and techniques used in the aforementioned studies, such as the use of pre-trained language models for the analysis of unstructured data.

\subsection{Role and Challenges of Shoutcasters}
\label{subsec:role-and-challenges-of-shoutcasters}

Shoutcasters play a centra role in both traditional sports and esports broadcasts, and therefore face several challenges in their profession that have been explored in the literature. Kempe-Cook et al. \cite{kempe2019Behind} conducted a study on amateur shoutcasters through a series of semi-structured interviews and participant observations. The study brought to attention the problem of amateur casters, who often start out alone, relying on reflection upon their own performance to improve their craft. 

\medskip 

Through the interviews, the study identifies two main types of casters, Play-by-Play and Color Commentary, depending on the focus of the commentary; narrating the live play, or providing analysis and insights, respectively. Experienced casters also provide \textit{design opportunities} for up-and-coming ones, both in and out of the game, such as social support for casting, improving in-game camerawork, and facilitating access to gaming data, in order to improve the quality of their broadcasts.

\medskip

Despite the length of the interviews, and the variety of participants, the study does not cover one of the main ways of audience engagement, which is through live chat interactions. This aspect of shoutcasting is often used in order to create a sense of excitement (also referred to as \textit{hype}) during the broadcast, and is present in both types of casters, and in League of Legends broadcasts. Omitting this aspect is a gap that needs to be filled, as there is a need to fully understand all aspects of shoutcasting, in order to be able to model it effectivelly, and possibly design AI systems to assist shoutcasters in their craft.

\medskip

Beyond esports, studies have developed tools to assist sports commentators during the broadcast, especially in the most difficult moments of the game. This is particularly relevant for games with multiple fast-paced events, both sports and videogames. In the case of the former, Lee et al. \cite{lee2025sportify}, designed a tool for basketball commentators. This tool aimed to aid with the task of processing all the difficult actions that happen during a basketball game, and helping the audience to understand them. An additional benefit was also to avoid the need of heavy visual aids, which often distract the audience from the game itself.

\medskip
% Might need a citation for the drama media, how anime characters talk to each other maybe.
The tool itself helped visualize tactical decisions, which represent a challenge to explain due to their technical level, as if the players were talking to each other, presenting the information in a more narrative way. The use of large language models was used to generate the commentary in real-time, and tests were conducted to evaluate the audience engagement and enjoyment of the tool. 

\medskip

A key insight obtained is another discursive function that was not covered in the prior study, yet still applies to esports, the importance of \textit{Storytelling} in commentary. The use of narrative elements to explain elements from the game enhanced user comprehension beyond just the technical aspects of the game, and it also increased audience enjoyment. However, the tests to verify this were limited in their scope, as they were conducted in a very controlled environment, and with a small sample of plays. As a part of the present study, we aim to improve the reach of tests, by using real broadcasts to obtain data, and including this function in the design of the discourse function labels. 

\medskip

In summary, the role of shoutcasters is acknowledged as crucial and therefore offers justification for the area of study. However, there are still several gaps in the understanding of the craft of shoutcasting, such as the role of live chat interactions, and the importance of storytelling in the commentary. These are important aspects to consider when modeling the structure and function of shoutcasting discourse, as they are key elements that contribute to audience engagement and enjoyment.

\subsection{Communicative Limitations of Shoutcasting}
\label{subsec:communicative-limitations-of-shoutcasting}

An alternative perspective on shoutcasting focuses on analyzing the communicative limitations of the craft, and use immersive technologies and AI tools to overcome them. 

\medskip 

The former approach was explored by Noma et al. \cite{noma2001new} in the domain of athletic events. The study aimed at overcoming the limitation of the audience not knowing how an athlete experiences physical strain during intense moments of competition. The method used was to design a Virtual Reality experience connected to the athlete's heart rate, allowing a user to experience a similar fatigue comparable to the athlete's by adjusting the speed of a treadmill. 

\medskip

The results of this study highlighted the potential of immersive technologies to enhance the audience experience, especially of those with lower skill levels. At the same time, this was one of the limitations of the study, as it seemed to only be ideal when working with ``ordinary'' users, and not more experienced ones. Despite this limitation, the study remains relevant, as more modern esports broadcasts incorporate Point of View (POV) perspectives, which allows the audience to further immerse themselves in the game, even if their skill level is not high. 

\medskip

A more recent and drastic approach to overcoming the communicative limitations of shoutcasting was through the use of AI tools to replace human shoutcasters. This was explored by Kim et al. \cite{kim2020automatic} in the context of baseball broadcasts. A pipeline of four Deep Learning models was designed to process a scene from a game and generate commentary through answering questions such as:

\begin{itemize}
  \item Who is doing what?
  \item In what area of the pitch is the action taking place?
  \item What can we expect to happen next?
\end{itemize}

\medskip

The results of such a system were mixed, as it was able to generate commentary that is relevant to the game, but still lacked proper domain knowledge to generate insightful commentary. As a result, the system had a high risk of delivering misleading information, requiring mitigation strategies such as multiple comments being done in a speculative manner. 

\medskip

Despite these limitations, the study is relevant for the present research, as it highlights the level of complexity that is required to generate high-quality commentary. This suggests that AI tools are better suited to assist human shoutcasters, rather than replace them. The pipeline of models used in the study can also provides guidelines for the design of AI systems, and generalized to the context of \textit{League of Legends} broadcasts.

\medskip

Overall, the communicative limitations of shoutcasting have been explored and attempts to overcome them have been developed. However, there remains a gap in the comprehensive understanding of shoutcasting as a communicative practice, particularly in esports, as well as in the design of AI tools to support shoutcasters. The present study aims to fill this gap by modeling the structure and function of shoutcasting discourse, and designing an AI system to assist shoutcasters in their craft, while also considering the communicative limitations of shoutcasting. 

\medskip

The following section explores in further detail the possibilities in discourse analysis, and the research that has been done in that domain, in order to find a suitable method for the analyzing shoutcasting discourse.